\clearpage\pagestyle{otMain}\fontsize{11}{15}\selectfont
\corrigesChapitre{13}{Tableaux croisés et fréquences}
\begin{blocLivre}{cor-13-1}
\begin{corrige}{1}{Lire les marges}
\demarche{Distinguer case intérieure et totaux marginaux.}
Le total général est 120. La marge de la ligne donne $n_A=60$ ; la marge de la colonne donne $n_C=60$. La case à leur croisement donne $n_{A\cap C}=36$. Ce dernier nombre ne représente ni tous les membres de $A$, ni tous les inscrits du club, mais ceux qui vérifient les deux conditions.
\end{corrige}
\end{blocLivre}
\begin{blocLivre}{cor-13-2}
\begin{corrige}{2}{Une fréquence conditionnelle}
\demarche{Restreindre le total à la ligne demandée.}
Parmi les 60 individus de $A$, 36 sont dans $C$ : la fréquence est $36/60=\res{60\,\%}$. Parmi les 60 de $B$, 24 sont dans $C$ : $24/60=\res{40\,\%}$. Chaque dénominateur est le total du groupe introduit par « parmi ».
\end{corrige}
\end{blocLivre}
\begin{blocLivre}{cor-13-3}
\begin{corrige}{3}{Inverser la référence}
\demarche{Conserver l’intersection et changer le dénominateur avec la question.}
La première fréquence vaut $30/50=\res{60\,\%}$ ; la seconde vaut $30/120=\res{25\,\%}$. Le numérateur reste identique parce qu’il décrit les mêmes 30 individus. Le dénominateur change avec la population conditionnée. Les deux résultats ne sont donc pas contradictoires.
\end{corrige}
\end{blocLivre}
\begin{blocLivre}{cor-13-4}
\begin{corrige}{4}{Compléter avec des pourcentages}
\demarche{Calculer les cases à partir des fréquences de leur groupe.}
Dans $A$, les satisfaits sont $0{,}75\times80=60$ et les autres 20. Dans $B$, ils sont $0{,}40\times120=48$ et les autres 72. Les totaux de colonnes sont $\res{108\text{ satisfaits et }92\text{ non satisfaits}}$. Le contrôle $60+20+48+72=200$ retrouve l’effectif général.
\end{corrige}
\end{blocLivre}
\begin{blocLivre}{cor-13-5}
\begin{corrige}{5}{Nominal ou ordinal ?}
\demarche{Chercher si un ordre des catégories a un sens.}
Le moyen de transport et la couleur des yeux sont des variables nominales : aucun ordre naturel unique n’est imposé. L’avis et le niveau sont ordinaux : les catégories ont un ordre explicite. Cela ne les transforme pas en nombres : attribuer des codes 1, 2, 3 ne prouve pas que les écarts entre catégories sont identiques.
\end{corrige}
\end{blocLivre}
\begin{blocLivre}{cor-13-6}
\begin{corrige}{6}{Marge et condition}
\demarche{Nommer le dénominateur avant d’effectuer le quotient.}
La fréquence globale du club est $60/120=\res{50\,\%}$. Parmi les 60 non-inscrits, 24 viennent du groupe $A$ : la seconde fréquence est $24/60=\res{40\,\%}$. Le premier calcul utilise le total général ; le deuxième utilise la marge de la colonne des non-inscrits.
\end{corrige}
\end{blocLivre}
\begin{blocLivre}{cor-13-7}
\begin{corrige}{7}{ET, OU et complémentaire}
\demarche{Utiliser l’inclusion-exclusion pour ne pas doubler l’intersection.}
« $A$ et $C$ » donne directement 36. Pour « $A$ ou $C$ », on additionne les totaux puis on retire l’intersection comptée deux fois : $60+60-36=\res{84}$. Enfin non $C$ contient $120-60=\res{60}$ personnes. Le « ou » mathématique inclut les personnes qui vérifient les deux conditions.
\label{fig-cor-13-7}
\figureGeo[corrige]{croise-venn}
\end{corrige}
\end{blocLivre}
\begin{blocLivre}{cor-13-8}
\begin{corrige}{8}{Du fichier au tableau}
\demarche{Compter les couples, pas seulement les deux variables séparées.}
On compte séparément chaque couple. Dans $A$, il y a 2 $O$ et 2 $N$ ; dans $B$, également 2 $O$ et 2 $N$. Chaque ligne et chaque colonne totalise 4 ; le total général est 8. Pour vérifier, barrer ou cocher chaque ligne du fichier après son comptage évite les oublis et les doubles comptes.
\end{corrige}
\end{blocLivre}
\begin{blocLivre}{cor-13-9}
\begin{corrige}{9}{Comparer des taux avec des effectifs}
\demarche{Comparer les rapports puis réunir les numérateurs et les effectifs.}
$A$ a un taux de $18/20=90\,\%$ ; $B$, $72/90=80\,\%$. Le taux de $A$ est plus élevé, même si $B$ compte plus de réussites en nombre. Globalement, il y a 90 réussites sur 110 personnes, soit $\res{9/11\approx81{,}8\,\%}$. Il faut distinguer effectifs et fréquences.
\end{corrige}
\end{blocLivre}
\begin{blocLivre}{cor-13-10}
\begin{corrige}{10}{Des données réelles de l’Insee}
\demarche{Lire une donnée réelle en conservant son champ et sa référence.}
Les trois dénominateurs sont différents. On calcule $5934049/32987021\approx\res{18{,}0\,\%}$, $5674817/35107259\approx\res{16{,}2\,\%}$ et $5674817/(5934049+5674817)\approx\res{48{,}9\,\%}$. La dernière référence est l’ensemble des moins de 15 ans, pas l’ensemble des femmes. Le champ, le millésime et le découpage géographique font partie de l’information.
\end{corrige}
\end{blocLivre}
\begin{blocLivre}{cor-13-11}
\begin{corrige}{11}{Un total à retrouver}
\demarche{Reconstituer le dénominateur manquant puis agréger les effectifs.}
Le groupe $A$ compte $24/0{,}60=\res{40}$ personnes. Dans l’ensemble, 39 personnes pratiquent l’activité parmi 90 : $39/90=\res{13/30\approx43{,}3\,\%}$. Il faut retrouver le total manquant avant de réunir les groupes ; les deux pourcentages ne se moyennent pas à poids égaux.
\end{corrige}
\end{blocLivre}
\begin{blocLivre}{cor-13-12}
\begin{corrige}{12}{Une annonce impossible}
\demarche{Vérifier la compatibilité d’une proportion avec un effectif entier.}
Exactement 30\,\% de 35 donnerait $35\times0{,}30=10{,}5$ personnes, impossible pour un effectif. L’annonce peut être cohérente comme pourcentage approximatif : 10 sur 35 font environ 28,6\,\%, et 11 sur 35 environ 31,4\,\%, tous deux proches de 30\,\% à une précision grossière. Il faut annoncer l’arrondi utilisé plutôt que le mot « exactement ».
\end{corrige}
\end{blocLivre}
\begin{blocLivre}{cor-13-13}
\begin{corrige}{13}{Des fréquences identiques}
\demarche{Factoriser une proportion commune avant la division.}
Si les effectifs sont $a$ et $b$, le nombre d’individus de $C$ est $0{,}4a+0{,}4b=0{,}4(a+b)$. En divisant par le total $a+b>0$, on obtient $\res{0{,}4=40\,\%}$. La taille relative des groupes n’a ici aucun effet parce que les deux fréquences conditionnelles sont égales.
\end{corrige}
\end{blocLivre}
\begin{blocLivre}{cor-13-14}
\begin{corrige}{14}{L’effet d’un recrutement}
\demarche{Faire évoluer simultanément le numérateur et le dénominateur.}
Après $n$ recrutements, leur part est $(30+n)/(80+n)$ avec $n\ge0$. Le dénominateur est positif : l’inégalité à $1/2$ équivaut à $60+2n\ge80+n$, soit $n\ge20$. Il faut $\res{20\text{ personnes}}$. Au seuil, la proportion vaut $50/100=1/2$. Le total augmente lui aussi lors du recrutement.
\end{corrige}
\end{blocLivre}
\begin{blocLivre}{cor-13-15}
\begin{corrige}{15}{Des groupes à pondérer}
\demarche{Calculer tous les effectifs globaux et interpréter l’effet des pondérations.}
Dans le premier groupe, $X$ a 90\,\% contre 80\,\% pour $Y$ ; dans le second, 60\,\% contre 50\,\%. Pourtant globalement, $X$ a $33/50=66\,\%$ et $Y$ $37/50=74\,\%$. Les répartitions des effectifs entre groupes sont différentes. Une comparaison globale ne conserve donc pas nécessairement l’ordre observé dans chaque groupe.
\end{corrige}
\end{blocLivre}
